% sec:rk-dormand-prince — The Dormand--Prince RK45 Pair
\section{The Dormand--Prince RK45 Pair}
\label{sec:rk-dormand-prince}

The classical RK4 method of the previous section advances the solution
with fourth-order accuracy, but provides no indication of how accurate
any given step actually was.  A step through a gentle region of the
right-hand side wastes function evaluations on accuracy that isn't
needed, while a step near the photon sphere --- where the geodesic
right-hand side varies rapidly --- may introduce errors far larger
than the tolerance demands.  Without an error estimate, the integrator
has two bad options: use a conservatively small step size everywhere
(slow), or use a fixed step size that risks corrupting the trajectory
near the photon sphere where geodesic deviation is exponential.  An
adaptive integrator needs a local error estimate at every step, and the
elegant solution is to embed a second method of lower order within the
same set of stages.

\paragraph{Embedded pairs.}

The idea is to equip the Butcher tableau with a second row of weights.
Given the $s$ stage values $k_1, \ldots, k_s$ already computed for the
primary method of order $p$, a companion method of order $\hat{p} < p$
constructs a second solution using different weights
$\hat{\mathbf{b}}$:
%
\begin{equation}
  \hat{y}_{n+1} = y_n + h \sum_{i=1}^{s} \hat{b}_i\,k_i \,.
  \label{eq:embedded-companion}
\end{equation}
%
Both solutions use the \emph{same} stages --- no additional evaluations
of $f$ are required.  The Butcher tableau notation extends naturally by
adding the companion weights as a second row below the line:
%
\begin{equation}
  \renewcommand{\arraystretch}{1.3}
  \begin{array}{c|c}
    \mathbf{c} & A \\[2pt]
    \hline
    & \mathbf{b}^{\!\top} \\
    & \hat{\mathbf{b}}^{\!\top}
  \end{array}
  \label{eq:embedded-tableau}
\end{equation}
%
The pair is denoted $p(\hat{p})$ --- for Dormand--Prince, $5(4)$ ---
meaning a fifth-order primary method with a fourth-order companion.
\histnote{Dormand and Prince published this pair in 1980
\cite{DormandPrince1980}; it has since become the default ODE
solver in most scientific computing libraries.}

The difference between the two solutions provides a local error
estimate.  Subtracting \cref{eq:embedded-companion} from the primary
update \cref{eq:rk-update} gives
%
\begin{equation}
  \delta y_{n+1} \;=\; y_{n+1} - \hat{y}_{n+1}
  \;=\; h \sum_{i=1}^{s} e_i\,k_i \,,
  \qquad e_i = b_i - \hat{b}_i \,.
  \label{eq:embedded-error}
\end{equation}
%
Because the primary solution has order $p = 5$ and the companion has
order $\hat{p} = 4$, this difference is $\delta y_{n+1} = \order{h^5}$:
it approximates the leading-order local truncation error of the
lower-order companion, up to $\order{h^6}$ corrections from the primary
method's own truncation error.  The integrator uses this estimate to
decide whether to accept the step and how to adjust the step size ---
machinery developed in \cref{sec:rk-adaptive}.

\paragraph{The Dormand--Prince 5(4) coefficients.}

Seven stages ($s = 7$) achieve fifth-order accuracy.  The Butcher barrier of \cref{eq:butcher-barrier}
demands at least six stages for order five; the purpose of the seventh
stage becomes clear in \cref{sec:rk-fsal}.  Dormand and Prince solved
the seventeen order conditions required for fifth-order accuracy and
used the remaining free parameters to minimise the leading
truncation-error coefficient of the fifth-order solution
\cite{DormandPrince1980}.  The resulting tableau is
\coderef{Geodesic}{Solver}
%
\begin{equation}
  \renewcommand{\arraystretch}{1.4}
  \footnotesize
  \begin{array}{c|ccccccc}
    0      & & & & & & & \\
    \tfrac{1}{5}  & \tfrac{1}{5} & & & & & & \\
    \tfrac{3}{10} & \tfrac{3}{40} & \tfrac{9}{40} & & & & & \\
    \tfrac{4}{5}  & \tfrac{44}{45} & -\tfrac{56}{15} & \tfrac{32}{9} & & & & \\
    \tfrac{8}{9}  & \tfrac{19372}{6561} & -\tfrac{25360}{2187} & \tfrac{64448}{6561} & -\tfrac{212}{729} & & & \\
    1      & \tfrac{9017}{3168} & -\tfrac{355}{33} & \tfrac{46732}{5247} & \tfrac{49}{176} & -\tfrac{5103}{18656} & & \\
    1      & \tfrac{35}{384} & 0 & \tfrac{500}{1113} & \tfrac{125}{192} & -\tfrac{2187}{6784} & \tfrac{11}{84} & \\[4pt]
    \hline
    \\[-4pt]
           & \tfrac{35}{384} & 0 & \tfrac{500}{1113} & \tfrac{125}{192} & -\tfrac{2187}{6784} & \tfrac{11}{84} & 0 \\
           & \tfrac{5179}{57600} & 0 & \tfrac{7571}{16695} & \tfrac{393}{640} & -\tfrac{92097}{339200} & \tfrac{187}{2100} & \tfrac{1}{40}
  \end{array}
  \label{eq:dp-tableau}
\end{equation}
%
Several features of this tableau deserve attention.  The nodes
$c_6 = c_7 = 1$ mean that stages six and seven both evaluate $f$ at
the far end of the step interval --- but at different $y$ values,
since they are constructed from different linear combinations of the
preceding stages.  The primary weights $\mathbf{b}$ (first row below
the line) are identical to the seventh row of $A$: the fifth-order
solution $y_{n+1}$ is precisely the point at which $k_7$ is evaluated.
This structural coincidence is the FSAL (First Same As Last) property
developed in \cref{sec:rk-fsal}.  The weight $b_2 = \hat{b}_2 = 0$
means that stage $k_2$ does not contribute directly to either solution,
though it influences the later stages through the coupling matrix.
\warnnote{Several coupling coefficients $a_{ij}$ are negative
--- for instance $a_{42} = -56/15$ and $a_{52} = -25360/2187$.  Negative
weights are typical of high-order methods: they allow cancellation of
higher-order error terms that positive coefficients alone cannot
achieve.}

\paragraph{The error coefficients.}

Computing $e_i = b_i - \hat{b}_i$ from \cref{eq:embedded-error}
gives the seven error coefficients used by the integrator:
%
\begin{equation}
  \renewcommand{\arraystretch}{1.5}
  \begin{array}{l|ccccccc}
    i & 1 & 2 & 3 & 4 & 5 & 6 & 7 \\
    \hline
    e_i & \tfrac{71}{57600} & 0 & -\tfrac{71}{16695} & \tfrac{71}{1920}
        & -\tfrac{17253}{339200} & \tfrac{22}{525} & -\tfrac{1}{40}
  \end{array}
  \label{eq:dp-error-coeffs}
\end{equation}
%
The codebase stores these pre-computed differences rather than the
companion weights $\hat{\mathbf{b}}$ themselves, since the error
estimate is the only place $\hat{\mathbf{b}}$ appears.  The most
consequential entry is $e_7 = -1/40$: because $b_7 = 0$ but
$\hat{b}_7 = 1/40$, the companion solution needs the seventh stage
evaluation even though the primary solution does not --- it is this
nonzero error coefficient, not just the FSAL reuse, that makes the
seventh stage indispensable.  The entry $e_2 = 0$ reflects the fact
that both methods assign zero weight to stage two; $k_2$
contributes to accuracy only indirectly, through the coupling matrix.
\implnote{The \texttt{dormandPrinceStep} function in
\codemodule{Geodesic.Solver} stores the error coefficients
\texttt{e1}\,\ldots\,\texttt{e7} directly and computes the error
vector as $h\sum e_i k_i$ in a single pass.}

\paragraph{Local extrapolation.}

The notation $5(4)$ indicates that the \emph{fifth-order} solution is
propagated --- used as the starting point for the next step --- while
the fourth-order companion serves only to estimate the error.  This
choice is called \emph{local extrapolation}, and at first glance it
seems backwards: shouldn't the integrator propagate the solution whose
error it can actually estimate?

The argument for local extrapolation is pragmatic.  The fifth-order
solution has local truncation error $\order{h^6}$, one order better
than the $\order{h^5}$ error of the companion, so propagating it
produces a more accurate trajectory.  The error estimate
$\delta y_{n+1} = \order{h^5}$ is conservative when applied to a
solution that is actually accurate to $\order{h^6}$: the step-size
controller typically accepts steps that are better than it believes,
and in practice rarely allows errors exceeding the tolerance.
\xrefnote{The alternative --- propagating the lower-order solution ---
is sometimes called \emph{local interpolation} and is used by some
methods optimised for dense output \cite{HairerNorsettWanner1993}.}

The error estimate $\delta y_{n+1}$ tells us how wrong a step was ---
but not yet what to do about it.  Two pieces of practical machinery
remain: reclaiming the cost of the seventh stage
(\cref{sec:rk-fsal}) and using $\delta y_{n+1}$ to choose the next
step size (\cref{sec:rk-adaptive}).
